%%%%%%%%%%%%%%%%%%%%%%%%%%%%%%%%%%%%%%%%%%%%%%%
%%%     REVIEW CRITERIA     %%%%
%%%
%%%     LDRD call boilerplate, typeset as front matter (roman page
%%%     numbers) ahead of the proposal.  Comment out the
%%%     %%%%%%%%%%%%%%%%%%%%%%%%%%%%%%%%%%%%%%%%%%%%%%%
%%%     REVIEW CRITERIA     %%%%
%%%
%%%     LDRD call boilerplate, typeset as front matter (roman page
%%%     numbers) ahead of the proposal.  Comment out the
%%%     %%%%%%%%%%%%%%%%%%%%%%%%%%%%%%%%%%%%%%%%%%%%%%%
%%%     REVIEW CRITERIA     %%%%
%%%
%%%     LDRD call boilerplate, typeset as front matter (roman page
%%%     numbers) ahead of the proposal.  Comment out the
%%%     %%%%%%%%%%%%%%%%%%%%%%%%%%%%%%%%%%%%%%%%%%%%%%%
%%%     REVIEW CRITERIA     %%%%
%%%
%%%     LDRD call boilerplate, typeset as front matter (roman page
%%%     numbers) ahead of the proposal.  Comment out the
%%%     \input{review_criteria} in the main file to drop it.
%%%%%%%%%%%%%%%%%%%%%%%%%%%%%%%%%%%%%%%%%%%%%%%

\pagenumbering{roman}
\begin{center}
    \large \textcolor{red}{\textbf{REVIEW CRITERIA}}
\end{center}
\vspace*{\baselineskip}

\textcolor{red}{\textbf{Problem, Hypothesis \& Technical Rationale:}}
\vspace*{0.5\baselineskip}

\textcolor{red}{\textit{(Suggested section length \(\sim \)1
    page)}}\\ \textcolor{red}{STATE YOUR INNOVATION. COMPARE TO STATE
  OF THE ART. DEFINE SUCCESS AND DETAIL YOUR OUTCOMES (Intellectual
  property? Preliminary data for another proposal? Fill a gap?
  Foundational understanding?)}\\ \textcolor{red}{COMMUNICATE “WHY
  NOW?”}
\vspace*{\baselineskip}

\textcolor{red}{Clearly articulate a technically sound concept that defines the problem well, employs appropriate methods or models, and incorporates testable elements aimed at producing early results or reducing key uncertainties.}
\vspace*{\baselineskip}

\textcolor{red}{Review criteria:}
\begin{quote}
\vspace{-0.3cm}
    \textcolor{red}{\textit{\textbf{Technical Vitality} review criteria weighting 40\%, covering \textbf{Clarity of Problem and Objectives} and \textbf{Innovation and Technical Merit.}}}
\end{quote}

\begin{itemize}[noitemsep,nolistsep]
    \item \textcolor{red}{\textit{What \textbf{technical problem or uncertainty} does this feasibility study address?}}
    \item \textcolor{red}{\textit{Why is resolving this uncertainty important?}}
    \item \textcolor{red}{\textit{What are the \textbf{specific hypotheses, assumptions, or risks} to be tested or reduced?}}
    \item \textcolor{red}{\textit{What is the \textbf{technical design} of the study?}}
    \item \textcolor{red}{\textit{What experiments, models, or analyses will be performed?}}
    \item \textcolor{red}{\textit{How do these tasks directly test the hypotheses or reduce risk?}}
    \item \textcolor{red}{\textit{What \textbf{metrics or criteria} will determine whether hypotheses are supported, disproved, or reframed?}}
    \item \textcolor{red}{\textit{How will the study capture \textbf{useful lessons}, even from negative or inconclusive results?}}
\end{itemize}
\vspace*{\baselineskip}

\textcolor{red}{\textbf{Mission Impact \& Institutional Agility:}}
\vspace*{0.5\baselineskip}

\textcolor{red}{\textit{(Suggested section length \(\sim \)$\frac{1}{2}$ page)}}\\
\textcolor{red}{Foster cutting edge science, technology, and engineering that ready the Laboratory to respond to current and future mission and national needs.}
\vspace*{\baselineskip}

\textcolor{red}{Review criteria:}\\
\begin{quote}
    \vspace{-0.4cm}\textcolor{red}{\textit{\textbf{Mission Agility} review criteria weighting 20\%, covering \textbf{Mission Agility} and \textbf{Mission Institutional Learning.}}}
\end{quote}

\begin{itemize}[noitemsep,nolistsep]
    \item \textcolor{red}{\textit{What is the potential \textbf{value of the insights} (positive or negative) for:}}
    \begin{itemize}[noitemsep,nolistsep]
        \item \textcolor{red}{\textit{the \textbf{technical community}}}
        \item \textcolor{red}{\textit{LANL's \textbf{mission space}, and}}
        \item \textcolor{red}{\textit{\textbf{future projects} that could build on these findings?}}
    \end{itemize}
    \item \textcolor{red}{\textit{How will the study contribute to \textbf{institutional learning} (roadmaps, datasets, methods, lessons-learned)?}}
\end{itemize}
\vspace*{\baselineskip}

\textcolor{red}{\textbf{Approach, Resources \& Path Forward:}}
\vspace*{0.5\baselineskip}

\textcolor{red}{\textit{(Suggested section length \(\sim \)$\frac{3}{4}$ page)}}\\
\textcolor{red}{Articulates a feasibility study that is well conceived, clearly organized, and appropriately scoped, with structure, methods, and milestones designed to produce testable outcomes and generate results that can guide future work.}
\vspace*{\baselineskip}

\textcolor{red}{Review criteria:}
\begin{quote}
\vspace{-0.4cm}
    \textcolor{red}{\textit{\textbf{Research Approach} review criteria total weight 40\%, covering \textbf{Testability and Risk Reduction, Clarity and Organization} and \textbf{Scope, Readiness, Budget, and Duration.}}}
\end{quote}

\begin{itemize}[noitemsep,nolistsep]
    \item \textcolor{red}{\textit{What is the \textbf{plan of work} for the study?}}
    \item \textcolor{red}{\textit{Key \textbf{{tasks, milestones, or deliverables.}}}}
    \item \textcolor{red}{\textit{How will your approach provide \textbf{evidence to validate, refine, or reject} the hypothesis?}}
    \item \textcolor{red}{\textit{What outputs will you produce (data, models, prototypes, analyses)?}}
    \item \textcolor{red}{\textit{What are the \textbf{success criteria}?}}
    \item \textcolor{red}{\textit{How will results - positive or negative - inform \textbf{next steps} (scale up, redesign, or close)?}}
    \item \textcolor{red}{\textit{What are the \textbf{non-technical risks} (facility access, staffing, regulatory, collaboration) and how will you mitigate them?}}
\end{itemize}
\vspace*{\baselineskip}

\textcolor{red}{Proposers should provide enough information for reviewers to assess the scope of work, budget and duration – and that the project is in compliance with \textbf{\href{https://ldrd.lanl.gov/q/public/policy?inline=true}{LDRD Rules and Procedures} (pdf)}}

\begin{itemize}[noitemsep,nolistsep]
    \item \textcolor{red}{\textit{Why are the \textbf{proposed budget} and \textbf{duration} appropriate? (optional bullets: Personnel, Equipment, Materials, Travel)}}
\end{itemize}
\vspace*{\baselineskip}

\textcolor{red}{Since the projects are such short time scales, it should be stated that the work can be turned on immediately if funded. If there is a new activity involved that needs approval, if special materials or lab equipment needs to be procured, if HPC accounts are not in place already, etc. this will raise questions for the reviewers on whether the proposed research can be completed in the time frame provided.}
\vspace*{\baselineskip}

\textcolor{red}{BUDGET RELATED:}\\
\textcolor{red}{If working with external collaborators, communicate timing, any preferred MOUs and academic partners. MAKE SURE YOU COMPLY WITH \textbf{\href{https://ldrd.lanl.gov/q/public/policy?inline=true}{LDRD Rules and Procedures} (pdf)}}\\
\textcolor{red}{FUNDING LIMITS AND PURCHASE DEADLINES.}
\begin{itemize}[noitemsep,nolistsep]
    \item \textcolor{red}{\textbf{Collaborators: External Collaborations are permitted.} No more than 25\% of the project budget may be spent at a collaborator institution \textbf{(7.3.5 Budget Share for Collaborators).}}
    \item \textcolor{red}{\textbf{Purchases:}}
    \begin{itemize}[noitemsep,nolistsep]
        \item \textcolor{red}{M\&S costs must not exceed 25\% of the project’s total budget \textbf{(7.3.2 Materials and Services (M\&S))}}
        \item \textcolor{red}{All single items \> \$5k must be received at least three-months before the project ends. \textbf{(7.4.2 Purchasing Limitations for Final-year Projects)}}
    \end{itemize}
\end{itemize}

\pagebreak
\pagenumbering{arabic}
\setcounter{page}{1}
 in the main file to drop it.
%%%%%%%%%%%%%%%%%%%%%%%%%%%%%%%%%%%%%%%%%%%%%%%

\pagenumbering{roman}
\begin{center}
    \large \textcolor{red}{\textbf{REVIEW CRITERIA}}
\end{center}
\vspace*{\baselineskip}

\textcolor{red}{\textbf{Problem, Hypothesis \& Technical Rationale:}}
\vspace*{0.5\baselineskip}

\textcolor{red}{\textit{(Suggested section length \(\sim \)1
    page)}}\\ \textcolor{red}{STATE YOUR INNOVATION. COMPARE TO STATE
  OF THE ART. DEFINE SUCCESS AND DETAIL YOUR OUTCOMES (Intellectual
  property? Preliminary data for another proposal? Fill a gap?
  Foundational understanding?)}\\ \textcolor{red}{COMMUNICATE “WHY
  NOW?”}
\vspace*{\baselineskip}

\textcolor{red}{Clearly articulate a technically sound concept that defines the problem well, employs appropriate methods or models, and incorporates testable elements aimed at producing early results or reducing key uncertainties.}
\vspace*{\baselineskip}

\textcolor{red}{Review criteria:}
\begin{quote}
\vspace{-0.3cm}
    \textcolor{red}{\textit{\textbf{Technical Vitality} review criteria weighting 40\%, covering \textbf{Clarity of Problem and Objectives} and \textbf{Innovation and Technical Merit.}}}
\end{quote}

\begin{itemize}[noitemsep,nolistsep]
    \item \textcolor{red}{\textit{What \textbf{technical problem or uncertainty} does this feasibility study address?}}
    \item \textcolor{red}{\textit{Why is resolving this uncertainty important?}}
    \item \textcolor{red}{\textit{What are the \textbf{specific hypotheses, assumptions, or risks} to be tested or reduced?}}
    \item \textcolor{red}{\textit{What is the \textbf{technical design} of the study?}}
    \item \textcolor{red}{\textit{What experiments, models, or analyses will be performed?}}
    \item \textcolor{red}{\textit{How do these tasks directly test the hypotheses or reduce risk?}}
    \item \textcolor{red}{\textit{What \textbf{metrics or criteria} will determine whether hypotheses are supported, disproved, or reframed?}}
    \item \textcolor{red}{\textit{How will the study capture \textbf{useful lessons}, even from negative or inconclusive results?}}
\end{itemize}
\vspace*{\baselineskip}

\textcolor{red}{\textbf{Mission Impact \& Institutional Agility:}}
\vspace*{0.5\baselineskip}

\textcolor{red}{\textit{(Suggested section length \(\sim \)$\frac{1}{2}$ page)}}\\
\textcolor{red}{Foster cutting edge science, technology, and engineering that ready the Laboratory to respond to current and future mission and national needs.}
\vspace*{\baselineskip}

\textcolor{red}{Review criteria:}\\
\begin{quote}
    \vspace{-0.4cm}\textcolor{red}{\textit{\textbf{Mission Agility} review criteria weighting 20\%, covering \textbf{Mission Agility} and \textbf{Mission Institutional Learning.}}}
\end{quote}

\begin{itemize}[noitemsep,nolistsep]
    \item \textcolor{red}{\textit{What is the potential \textbf{value of the insights} (positive or negative) for:}}
    \begin{itemize}[noitemsep,nolistsep]
        \item \textcolor{red}{\textit{the \textbf{technical community}}}
        \item \textcolor{red}{\textit{LANL's \textbf{mission space}, and}}
        \item \textcolor{red}{\textit{\textbf{future projects} that could build on these findings?}}
    \end{itemize}
    \item \textcolor{red}{\textit{How will the study contribute to \textbf{institutional learning} (roadmaps, datasets, methods, lessons-learned)?}}
\end{itemize}
\vspace*{\baselineskip}

\textcolor{red}{\textbf{Approach, Resources \& Path Forward:}}
\vspace*{0.5\baselineskip}

\textcolor{red}{\textit{(Suggested section length \(\sim \)$\frac{3}{4}$ page)}}\\
\textcolor{red}{Articulates a feasibility study that is well conceived, clearly organized, and appropriately scoped, with structure, methods, and milestones designed to produce testable outcomes and generate results that can guide future work.}
\vspace*{\baselineskip}

\textcolor{red}{Review criteria:}
\begin{quote}
\vspace{-0.4cm}
    \textcolor{red}{\textit{\textbf{Research Approach} review criteria total weight 40\%, covering \textbf{Testability and Risk Reduction, Clarity and Organization} and \textbf{Scope, Readiness, Budget, and Duration.}}}
\end{quote}

\begin{itemize}[noitemsep,nolistsep]
    \item \textcolor{red}{\textit{What is the \textbf{plan of work} for the study?}}
    \item \textcolor{red}{\textit{Key \textbf{{tasks, milestones, or deliverables.}}}}
    \item \textcolor{red}{\textit{How will your approach provide \textbf{evidence to validate, refine, or reject} the hypothesis?}}
    \item \textcolor{red}{\textit{What outputs will you produce (data, models, prototypes, analyses)?}}
    \item \textcolor{red}{\textit{What are the \textbf{success criteria}?}}
    \item \textcolor{red}{\textit{How will results - positive or negative - inform \textbf{next steps} (scale up, redesign, or close)?}}
    \item \textcolor{red}{\textit{What are the \textbf{non-technical risks} (facility access, staffing, regulatory, collaboration) and how will you mitigate them?}}
\end{itemize}
\vspace*{\baselineskip}

\textcolor{red}{Proposers should provide enough information for reviewers to assess the scope of work, budget and duration – and that the project is in compliance with \textbf{\href{https://ldrd.lanl.gov/q/public/policy?inline=true}{LDRD Rules and Procedures} (pdf)}}

\begin{itemize}[noitemsep,nolistsep]
    \item \textcolor{red}{\textit{Why are the \textbf{proposed budget} and \textbf{duration} appropriate? (optional bullets: Personnel, Equipment, Materials, Travel)}}
\end{itemize}
\vspace*{\baselineskip}

\textcolor{red}{Since the projects are such short time scales, it should be stated that the work can be turned on immediately if funded. If there is a new activity involved that needs approval, if special materials or lab equipment needs to be procured, if HPC accounts are not in place already, etc. this will raise questions for the reviewers on whether the proposed research can be completed in the time frame provided.}
\vspace*{\baselineskip}

\textcolor{red}{BUDGET RELATED:}\\
\textcolor{red}{If working with external collaborators, communicate timing, any preferred MOUs and academic partners. MAKE SURE YOU COMPLY WITH \textbf{\href{https://ldrd.lanl.gov/q/public/policy?inline=true}{LDRD Rules and Procedures} (pdf)}}\\
\textcolor{red}{FUNDING LIMITS AND PURCHASE DEADLINES.}
\begin{itemize}[noitemsep,nolistsep]
    \item \textcolor{red}{\textbf{Collaborators: External Collaborations are permitted.} No more than 25\% of the project budget may be spent at a collaborator institution \textbf{(7.3.5 Budget Share for Collaborators).}}
    \item \textcolor{red}{\textbf{Purchases:}}
    \begin{itemize}[noitemsep,nolistsep]
        \item \textcolor{red}{M\&S costs must not exceed 25\% of the project’s total budget \textbf{(7.3.2 Materials and Services (M\&S))}}
        \item \textcolor{red}{All single items \> \$5k must be received at least three-months before the project ends. \textbf{(7.4.2 Purchasing Limitations for Final-year Projects)}}
    \end{itemize}
\end{itemize}

\pagebreak
\pagenumbering{arabic}
\setcounter{page}{1}
 in the main file to drop it.
%%%%%%%%%%%%%%%%%%%%%%%%%%%%%%%%%%%%%%%%%%%%%%%

\pagenumbering{roman}
\begin{center}
    \large \textcolor{red}{\textbf{REVIEW CRITERIA}}
\end{center}
\vspace*{\baselineskip}

\textcolor{red}{\textbf{Problem, Hypothesis \& Technical Rationale:}}
\vspace*{0.5\baselineskip}

\textcolor{red}{\textit{(Suggested section length \(\sim \)1
    page)}}\\ \textcolor{red}{STATE YOUR INNOVATION. COMPARE TO STATE
  OF THE ART. DEFINE SUCCESS AND DETAIL YOUR OUTCOMES (Intellectual
  property? Preliminary data for another proposal? Fill a gap?
  Foundational understanding?)}\\ \textcolor{red}{COMMUNICATE “WHY
  NOW?”}
\vspace*{\baselineskip}

\textcolor{red}{Clearly articulate a technically sound concept that defines the problem well, employs appropriate methods or models, and incorporates testable elements aimed at producing early results or reducing key uncertainties.}
\vspace*{\baselineskip}

\textcolor{red}{Review criteria:}
\begin{quote}
\vspace{-0.3cm}
    \textcolor{red}{\textit{\textbf{Technical Vitality} review criteria weighting 40\%, covering \textbf{Clarity of Problem and Objectives} and \textbf{Innovation and Technical Merit.}}}
\end{quote}

\begin{itemize}[noitemsep,nolistsep]
    \item \textcolor{red}{\textit{What \textbf{technical problem or uncertainty} does this feasibility study address?}}
    \item \textcolor{red}{\textit{Why is resolving this uncertainty important?}}
    \item \textcolor{red}{\textit{What are the \textbf{specific hypotheses, assumptions, or risks} to be tested or reduced?}}
    \item \textcolor{red}{\textit{What is the \textbf{technical design} of the study?}}
    \item \textcolor{red}{\textit{What experiments, models, or analyses will be performed?}}
    \item \textcolor{red}{\textit{How do these tasks directly test the hypotheses or reduce risk?}}
    \item \textcolor{red}{\textit{What \textbf{metrics or criteria} will determine whether hypotheses are supported, disproved, or reframed?}}
    \item \textcolor{red}{\textit{How will the study capture \textbf{useful lessons}, even from negative or inconclusive results?}}
\end{itemize}
\vspace*{\baselineskip}

\textcolor{red}{\textbf{Mission Impact \& Institutional Agility:}}
\vspace*{0.5\baselineskip}

\textcolor{red}{\textit{(Suggested section length \(\sim \)$\frac{1}{2}$ page)}}\\
\textcolor{red}{Foster cutting edge science, technology, and engineering that ready the Laboratory to respond to current and future mission and national needs.}
\vspace*{\baselineskip}

\textcolor{red}{Review criteria:}\\
\begin{quote}
    \vspace{-0.4cm}\textcolor{red}{\textit{\textbf{Mission Agility} review criteria weighting 20\%, covering \textbf{Mission Agility} and \textbf{Mission Institutional Learning.}}}
\end{quote}

\begin{itemize}[noitemsep,nolistsep]
    \item \textcolor{red}{\textit{What is the potential \textbf{value of the insights} (positive or negative) for:}}
    \begin{itemize}[noitemsep,nolistsep]
        \item \textcolor{red}{\textit{the \textbf{technical community}}}
        \item \textcolor{red}{\textit{LANL's \textbf{mission space}, and}}
        \item \textcolor{red}{\textit{\textbf{future projects} that could build on these findings?}}
    \end{itemize}
    \item \textcolor{red}{\textit{How will the study contribute to \textbf{institutional learning} (roadmaps, datasets, methods, lessons-learned)?}}
\end{itemize}
\vspace*{\baselineskip}

\textcolor{red}{\textbf{Approach, Resources \& Path Forward:}}
\vspace*{0.5\baselineskip}

\textcolor{red}{\textit{(Suggested section length \(\sim \)$\frac{3}{4}$ page)}}\\
\textcolor{red}{Articulates a feasibility study that is well conceived, clearly organized, and appropriately scoped, with structure, methods, and milestones designed to produce testable outcomes and generate results that can guide future work.}
\vspace*{\baselineskip}

\textcolor{red}{Review criteria:}
\begin{quote}
\vspace{-0.4cm}
    \textcolor{red}{\textit{\textbf{Research Approach} review criteria total weight 40\%, covering \textbf{Testability and Risk Reduction, Clarity and Organization} and \textbf{Scope, Readiness, Budget, and Duration.}}}
\end{quote}

\begin{itemize}[noitemsep,nolistsep]
    \item \textcolor{red}{\textit{What is the \textbf{plan of work} for the study?}}
    \item \textcolor{red}{\textit{Key \textbf{{tasks, milestones, or deliverables.}}}}
    \item \textcolor{red}{\textit{How will your approach provide \textbf{evidence to validate, refine, or reject} the hypothesis?}}
    \item \textcolor{red}{\textit{What outputs will you produce (data, models, prototypes, analyses)?}}
    \item \textcolor{red}{\textit{What are the \textbf{success criteria}?}}
    \item \textcolor{red}{\textit{How will results - positive or negative - inform \textbf{next steps} (scale up, redesign, or close)?}}
    \item \textcolor{red}{\textit{What are the \textbf{non-technical risks} (facility access, staffing, regulatory, collaboration) and how will you mitigate them?}}
\end{itemize}
\vspace*{\baselineskip}

\textcolor{red}{Proposers should provide enough information for reviewers to assess the scope of work, budget and duration – and that the project is in compliance with \textbf{\href{https://ldrd.lanl.gov/q/public/policy?inline=true}{LDRD Rules and Procedures} (pdf)}}

\begin{itemize}[noitemsep,nolistsep]
    \item \textcolor{red}{\textit{Why are the \textbf{proposed budget} and \textbf{duration} appropriate? (optional bullets: Personnel, Equipment, Materials, Travel)}}
\end{itemize}
\vspace*{\baselineskip}

\textcolor{red}{Since the projects are such short time scales, it should be stated that the work can be turned on immediately if funded. If there is a new activity involved that needs approval, if special materials or lab equipment needs to be procured, if HPC accounts are not in place already, etc. this will raise questions for the reviewers on whether the proposed research can be completed in the time frame provided.}
\vspace*{\baselineskip}

\textcolor{red}{BUDGET RELATED:}\\
\textcolor{red}{If working with external collaborators, communicate timing, any preferred MOUs and academic partners. MAKE SURE YOU COMPLY WITH \textbf{\href{https://ldrd.lanl.gov/q/public/policy?inline=true}{LDRD Rules and Procedures} (pdf)}}\\
\textcolor{red}{FUNDING LIMITS AND PURCHASE DEADLINES.}
\begin{itemize}[noitemsep,nolistsep]
    \item \textcolor{red}{\textbf{Collaborators: External Collaborations are permitted.} No more than 25\% of the project budget may be spent at a collaborator institution \textbf{(7.3.5 Budget Share for Collaborators).}}
    \item \textcolor{red}{\textbf{Purchases:}}
    \begin{itemize}[noitemsep,nolistsep]
        \item \textcolor{red}{M\&S costs must not exceed 25\% of the project’s total budget \textbf{(7.3.2 Materials and Services (M\&S))}}
        \item \textcolor{red}{All single items \> \$5k must be received at least three-months before the project ends. \textbf{(7.4.2 Purchasing Limitations for Final-year Projects)}}
    \end{itemize}
\end{itemize}

\pagebreak
\pagenumbering{arabic}
\setcounter{page}{1}
 in the main file to drop it.
%%%%%%%%%%%%%%%%%%%%%%%%%%%%%%%%%%%%%%%%%%%%%%%

\pagenumbering{roman}
\begin{center}
    \large \textcolor{red}{\textbf{REVIEW CRITERIA}}
\end{center}
\vspace*{\baselineskip}

\textcolor{red}{\textbf{Problem, Hypothesis \& Technical Rationale:}}
\vspace*{0.5\baselineskip}

\textcolor{red}{\textit{(Suggested section length \(\sim \)1
    page)}}\\ \textcolor{red}{STATE YOUR INNOVATION. COMPARE TO STATE
  OF THE ART. DEFINE SUCCESS AND DETAIL YOUR OUTCOMES (Intellectual
  property? Preliminary data for another proposal? Fill a gap?
  Foundational understanding?)}\\ \textcolor{red}{COMMUNICATE “WHY
  NOW?”}
\vspace*{\baselineskip}

\textcolor{red}{Clearly articulate a technically sound concept that defines the problem well, employs appropriate methods or models, and incorporates testable elements aimed at producing early results or reducing key uncertainties.}
\vspace*{\baselineskip}

\textcolor{red}{Review criteria:}
\begin{quote}
\vspace{-0.3cm}
    \textcolor{red}{\textit{\textbf{Technical Vitality} review criteria weighting 40\%, covering \textbf{Clarity of Problem and Objectives} and \textbf{Innovation and Technical Merit.}}}
\end{quote}

\begin{itemize}[noitemsep,nolistsep]
    \item \textcolor{red}{\textit{What \textbf{technical problem or uncertainty} does this feasibility study address?}}
    \item \textcolor{red}{\textit{Why is resolving this uncertainty important?}}
    \item \textcolor{red}{\textit{What are the \textbf{specific hypotheses, assumptions, or risks} to be tested or reduced?}}
    \item \textcolor{red}{\textit{What is the \textbf{technical design} of the study?}}
    \item \textcolor{red}{\textit{What experiments, models, or analyses will be performed?}}
    \item \textcolor{red}{\textit{How do these tasks directly test the hypotheses or reduce risk?}}
    \item \textcolor{red}{\textit{What \textbf{metrics or criteria} will determine whether hypotheses are supported, disproved, or reframed?}}
    \item \textcolor{red}{\textit{How will the study capture \textbf{useful lessons}, even from negative or inconclusive results?}}
\end{itemize}
\vspace*{\baselineskip}

\textcolor{red}{\textbf{Mission Impact \& Institutional Agility:}}
\vspace*{0.5\baselineskip}

\textcolor{red}{\textit{(Suggested section length \(\sim \)$\frac{1}{2}$ page)}}\\
\textcolor{red}{Foster cutting edge science, technology, and engineering that ready the Laboratory to respond to current and future mission and national needs.}
\vspace*{\baselineskip}

\textcolor{red}{Review criteria:}\\
\begin{quote}
    \vspace{-0.4cm}\textcolor{red}{\textit{\textbf{Mission Agility} review criteria weighting 20\%, covering \textbf{Mission Agility} and \textbf{Mission Institutional Learning.}}}
\end{quote}

\begin{itemize}[noitemsep,nolistsep]
    \item \textcolor{red}{\textit{What is the potential \textbf{value of the insights} (positive or negative) for:}}
    \begin{itemize}[noitemsep,nolistsep]
        \item \textcolor{red}{\textit{the \textbf{technical community}}}
        \item \textcolor{red}{\textit{LANL's \textbf{mission space}, and}}
        \item \textcolor{red}{\textit{\textbf{future projects} that could build on these findings?}}
    \end{itemize}
    \item \textcolor{red}{\textit{How will the study contribute to \textbf{institutional learning} (roadmaps, datasets, methods, lessons-learned)?}}
\end{itemize}
\vspace*{\baselineskip}

\textcolor{red}{\textbf{Approach, Resources \& Path Forward:}}
\vspace*{0.5\baselineskip}

\textcolor{red}{\textit{(Suggested section length \(\sim \)$\frac{3}{4}$ page)}}\\
\textcolor{red}{Articulates a feasibility study that is well conceived, clearly organized, and appropriately scoped, with structure, methods, and milestones designed to produce testable outcomes and generate results that can guide future work.}
\vspace*{\baselineskip}

\textcolor{red}{Review criteria:}
\begin{quote}
\vspace{-0.4cm}
    \textcolor{red}{\textit{\textbf{Research Approach} review criteria total weight 40\%, covering \textbf{Testability and Risk Reduction, Clarity and Organization} and \textbf{Scope, Readiness, Budget, and Duration.}}}
\end{quote}

\begin{itemize}[noitemsep,nolistsep]
    \item \textcolor{red}{\textit{What is the \textbf{plan of work} for the study?}}
    \item \textcolor{red}{\textit{Key \textbf{{tasks, milestones, or deliverables.}}}}
    \item \textcolor{red}{\textit{How will your approach provide \textbf{evidence to validate, refine, or reject} the hypothesis?}}
    \item \textcolor{red}{\textit{What outputs will you produce (data, models, prototypes, analyses)?}}
    \item \textcolor{red}{\textit{What are the \textbf{success criteria}?}}
    \item \textcolor{red}{\textit{How will results - positive or negative - inform \textbf{next steps} (scale up, redesign, or close)?}}
    \item \textcolor{red}{\textit{What are the \textbf{non-technical risks} (facility access, staffing, regulatory, collaboration) and how will you mitigate them?}}
\end{itemize}
\vspace*{\baselineskip}

\textcolor{red}{Proposers should provide enough information for reviewers to assess the scope of work, budget and duration – and that the project is in compliance with \textbf{\href{https://ldrd.lanl.gov/q/public/policy?inline=true}{LDRD Rules and Procedures} (pdf)}}

\begin{itemize}[noitemsep,nolistsep]
    \item \textcolor{red}{\textit{Why are the \textbf{proposed budget} and \textbf{duration} appropriate? (optional bullets: Personnel, Equipment, Materials, Travel)}}
\end{itemize}
\vspace*{\baselineskip}

\textcolor{red}{Since the projects are such short time scales, it should be stated that the work can be turned on immediately if funded. If there is a new activity involved that needs approval, if special materials or lab equipment needs to be procured, if HPC accounts are not in place already, etc. this will raise questions for the reviewers on whether the proposed research can be completed in the time frame provided.}
\vspace*{\baselineskip}

\textcolor{red}{BUDGET RELATED:}\\
\textcolor{red}{If working with external collaborators, communicate timing, any preferred MOUs and academic partners. MAKE SURE YOU COMPLY WITH \textbf{\href{https://ldrd.lanl.gov/q/public/policy?inline=true}{LDRD Rules and Procedures} (pdf)}}\\
\textcolor{red}{FUNDING LIMITS AND PURCHASE DEADLINES.}
\begin{itemize}[noitemsep,nolistsep]
    \item \textcolor{red}{\textbf{Collaborators: External Collaborations are permitted.} No more than 25\% of the project budget may be spent at a collaborator institution \textbf{(7.3.5 Budget Share for Collaborators).}}
    \item \textcolor{red}{\textbf{Purchases:}}
    \begin{itemize}[noitemsep,nolistsep]
        \item \textcolor{red}{M\&S costs must not exceed 25\% of the project’s total budget \textbf{(7.3.2 Materials and Services (M\&S))}}
        \item \textcolor{red}{All single items \> \$5k must be received at least three-months before the project ends. \textbf{(7.4.2 Purchasing Limitations for Final-year Projects)}}
    \end{itemize}
\end{itemize}

\pagebreak
\pagenumbering{arabic}
\setcounter{page}{1}
